(a) The projections induce a natural map $p_1^*\Omega_{X/S}\oplus p_2^*\Omega_{Y/S}\to\Omega_{X\times_S Y/S}$ by (8.11). On affine charts $\operatorname{Spec}B$ and $\operatorname{Spec}C$ over $\operatorname{Spec}A$, its inverse is induced by the derivation
\[B\otimes_A C\to(\Omega_{B/A}\otimes_A C)\oplus(B\otimes_A\Omega_{C/A}),\qquad b\otimes c\mapsto(db\otimes c,b\otimes dc).\]
These maps are inverse on the generators $d(b\otimes c)$ and commute with restriction. Together with (8.10), this gives the required isomorphism.

(b) The product is nonsingular: at a closed point, (a) gives cotangent dimension $\dim X+\dim Y=\dim(X\times Y)$, so (8.7) gives regularity; all other local rings are localizations of these. Taking the highest exterior power of (a) gives $\omega_{X\times Y}\cong p_1^*\omega_X\otimes p_2^*\omega_Y$.

(c) By (8.20), $\omega_Y\cong\mathcal O_Y(3-3)\cong\mathcal O_Y$. Thus (b) gives $\omega_X\cong\mathcal O_X$. Since $X$ is an integral projective variety, $\Gamma(X,\mathcal O_X)=k$, and $p_g(X)=1$. On the other hand, \exref{I.7.2}(b),(e) gives $p_a(Y)=1$ and $p_a(X)=p_a(Y)^2-2p_a(Y)=-1$.
